\documentclass[11pt]{article}

\usepackage[margin=1in]{geometry}
\usepackage{amsmath,amssymb,amsthm}
\usepackage{graphicx}
\usepackage[colorlinks=true,linkcolor=blue,citecolor=blue,urlcolor=blue]{hyperref}

\newtheorem{theorem}{Theorem}
\newtheorem{lemma}{Lemma}
\newtheorem{remark}{Remark}

\newcommand{\R}{\mathbb{R}}
\newcommand{\E}{\mathbb{E}}
\newcommand{\Prob}{\mathbb{P}}
\newcommand{\norm}[1]{\left\lVert #1\right\rVert}
\newcommand{\abs}[1]{\left\lvert #1\right\rvert}
\newcommand{\eps}{\varepsilon}

\title{A p-stable solution of Casey's metric-preserving dimension-reduction conjecture above a fixed scale}
\author{Agent lane 16, run fa\_banach\_001}
\date{July 3, 2026}

\begin{document}
\maketitle

\section{Source and target}

The source paper is Michael P. Casey, ``Linear Dimension Reduction Approximately
Preserving a Function of the 1-Norm'', arXiv:1906.03536, Electronic Journal of
Statistics.  On page 3, after proving the Cauchy case, the paper states the
following conjecture for p-stable projections.

\begin{quote}
For \(1<p<2\), replace the entries \(F_{ij}\) by iid standard p-stable random
variables and set \(\mu(\lambda)=\E \xi(\lambda F_{11})\).  With \(F\) and
\(\mu\) modified as above, and \(\rho,\epsilon,k\) as in Theorem 1.0.1, the
bound
\[
\mu\left(\frac{\norm{x-y}_p}{1+\epsilon}\right)
\leq \rho(Fx,Fy)
\leq \mu\bigl((1+\epsilon)\norm{x-y}_p\bigr)
\]
holds for all \(x,y\in P\) with \(\norm{x-y}_p=\Omega(1)\).
\end{quote}

The paper defines
\[
  \rho(u,v)=\frac1k\sum_{i=1}^k \xi(\abs{u_i-v_i}),\qquad
  \xi(t)=\log(1+\sqrt t)+\frac12\log(1+t),\quad t\geq 0.
\]
Since \(\xi\) is defined on nonnegative arguments and the proof throughout uses
\(\xi(\lambda\abs{W})\), the statement below uses the consistent interpretation
\[
  \mu_p(\lambda)=\E\,\xi(\lambda \abs{W}),
\]
where \(W\) is a standard symmetric p-stable random variable.

\begin{center}
\includegraphics[width=0.95\linewidth]{figures/open_problem_crop.png}
\end{center}

\section{Result}

\begin{theorem}[Fixed-threshold p-stable form of the conjecture]
\label{thm:main}
Fix \(1<p<2\), \(\lambda_0>0\), and \(a>0\).  Let \(P\subset\R^D\) be an
\(N\)-point set, \(N\geq2\), and let \(F:\R^D\to\R^k\) have iid standard
symmetric p-stable entries, normalized so that
\(\E e^{itW}=e^{-\abs{t}^p}\).  Define \(\rho\), \(\xi\), and
\(\mu_p(\lambda)=\E\xi(\lambda\abs W)\) as above.

There is a constant \(C=C(p,\lambda_0,a)\) such that, for every \(0<\eps<1\),
if
\[
   k\geq
   \frac{C}{\eps^2(1-\eps)^2}\log N,
\]
then with probability at least \(1-N^{-a}\), every pair \(x,y\in P\) with
\(\norm{x-y}_p\geq\lambda_0\) satisfies
\[
\mu_p\left(\frac{\norm{x-y}_p}{1+\eps}\right)
\leq \rho(Fx,Fy)
\leq \mu_p\bigl((1+\eps)\norm{x-y}_p\bigr).
\]
\end{theorem}

Thus the conjecture is proved under the standard interpretation of
\(\norm{x-y}_p=\Omega(1)\) as a fixed positive lower bound.  The proof does not
claim uniformity when that hidden lower bound is allowed to tend to zero with
\(N\), \(\eps\), or \(p\).

\section{Proof intuition}

The Cauchy proof in the source paper computes first and second moments of
\(\xi(\lambda\abs X)\) explicitly.  For general p-stable \(W\), those explicit
formulas are not available, but the special form of \(\xi\) gives a simpler
uniform route.  Put
\[
   \phi(s)=\xi(e^s)
      =\log(1+e^{s/2})+\frac12\log(1+e^s).
\]
Then \(0\leq\phi'(s)\leq1\), so changing \(\lambda\) only translates a
1-Lipschitz function of \(S=\log\abs W\).  The variable \(S\) has exponential
tails in both directions: positive tails follow from finite \(q\)-moments of
p-stable variables, and negative tails follow from boundedness of the stable
density near zero.  Therefore the centered variables
\(\xi(\lambda\abs W)-\mu_p(\lambda)\) obey a Bernstein inequality uniformly in
\(\lambda\).

The remaining point is deterministic.  Above a fixed scale \(\lambda_0\), the
increase of \(\mu_p\) from \(\lambda/(1+\eps)\) to \(\lambda\), and from
\(\lambda\) to \((1+\eps)\lambda\), is at least a positive constant times
\(\eps\).  Uniform concentration within that gap for all pairs proves the
conjectured inequalities.

\section{Proof}

\begin{lemma}[Uniform sub-exponential concentration]
\label{lem:concentration}
For every \(1<p<2\) there are constants \(K_p,c_p>0\) such that, for all
\(\lambda>0\), all \(0<t\leq1\), and iid copies \(W_1,\ldots,W_k\) of \(W\),
\[
\Prob\left\{
\left|\frac1k\sum_{j=1}^k
\bigl(\xi(\lambda\abs{W_j})-\mu_p(\lambda)\bigr)\right|>t
\right\}
\leq 2\exp(-c_p k t^2).
\]
\end{lemma}

\begin{proof}
Let \(S=\log\abs W\) and \(\phi(s)=\xi(e^s)\).  Since
\[
 \phi'(s)=
 \frac12\frac{e^{s/2}}{1+e^{s/2}}
 +\frac12\frac{e^s}{1+e^s},
\]
\(\phi\) is 1-Lipschitz.  For \(b=\log\lambda\),
\[
\abs{\phi(b+S)-\E\phi(b+S)}
\leq \abs{\phi(b+S)-\phi(b)}
     +\E\abs{\phi(b+S)-\phi(b)}
\leq \abs S+\E\abs S .
\]
Choose any \(0<q<1\).  A standard symmetric p-stable variable has
\(\E\abs W^q<\infty\), since \(q<p\), and has a bounded continuous density near
zero, hence \(\E\abs W^{-q}<\infty\).  Consequently
\[
  \E e^{q\abs S}=\E\max(\abs W^q,\abs W^{-q})<\infty .
\]
Thus the centered variables
\(\phi(b+S)-\E\phi(b+S)\) have sub-exponential norm bounded by a constant
depending only on \(p\), uniformly over \(b\in\R\).  Bernstein's inequality for
averages of centered sub-exponential variables gives
\[
\Prob\left\{\left|\frac1k\sum_{j=1}^k
\bigl(\xi(\lambda\abs{W_j})-\mu_p(\lambda)\bigr)\right|>t\right\}
\leq
2\exp\{-c k\min(t^2/K_p^2,t/K_p)\}.
\]
Absorbing constants and using \(0<t\leq1\) proves the lemma.
\end{proof}

\begin{lemma}[A fixed-scale gap for \(\mu_p\)]
\label{lem:gap}
For every \(1<p<2\) and \(\lambda_0>0\), there is
\(\gamma=\gamma(p,\lambda_0)>0\) such that, for all
\(\lambda\geq\lambda_0\) and \(0<\eps<1\),
\[
 \mu_p((1+\eps)\lambda)-\mu_p(\lambda)\geq \gamma\eps
\quad\text{and}\quad
 \mu_p(\lambda)-\mu_p\!\left(\frac{\lambda}{1+\eps}\right)\geq \gamma\eps .
\]
\end{lemma}

\begin{proof}
For \(t>0\), set
\[
  h(t)=t\xi'(t)
  =\frac12\frac{\sqrt t}{1+\sqrt t}
   +\frac12\frac{t}{1+t}.
\]
The function \(h\) is increasing, bounded, and positive on \((0,\infty)\).
For \(r>0\) and \(A>B>0\),
\[
  \xi(Ar)-\xi(Br)=\int_B^A h(ur)\,\frac{du}{u}.
\]
By boundedness of \(h\), Fubini applies.  With \(A=1+\eps\),
\[
\mu_p(A\lambda)-\mu_p(\lambda)
=\int_1^A \E h(u\lambda\abs W)\,\frac{du}{u}
\geq \log A\,\E h(\lambda_0\abs W/2),
\]
where the last lower bound is deliberately weakened for later symmetry.
Likewise,
\[
\mu_p(\lambda)-\mu_p(\lambda/A)
=\int_{1/A}^1 \E h(u\lambda\abs W)\,\frac{du}{u}
\geq \log A\,\E h(\lambda_0\abs W/2),
\]
because \(A<2\), \(u\geq 1/A\), and \(\lambda\geq\lambda_0\).  Since
\(\log(1+\eps)\geq \eps/2\) for \(0<\eps<1\), the result follows with
\[
  \gamma=\frac12\,\E h(\lambda_0\abs W/2)>0 .
\]
\end{proof}

\begin{proof}[Proof of Theorem~\ref{thm:main}]
Fix a pair \(x,y\in P\), put \(v=x-y\), and let
\(\lambda=\norm v_p\).  By p-stability, each coordinate of \(Fv\) has the same
law as \(\lambda W\), independently over the \(k\) coordinates.  Hence
\[
  \rho(Fx,Fy)
  \stackrel{d}{=}
  \frac1k\sum_{j=1}^k \xi(\lambda\abs{W_j}).
\]
Assume \(\lambda\geq\lambda_0\).  Lemma~\ref{lem:gap} gives a gap
\(\gamma\eps\) from \(\mu_p(\lambda)\) to each of the two desired boundary
values.  Therefore the event
\[
 \left|\rho(Fx,Fy)-\mu_p(\lambda)\right|\leq\gamma\eps
\]
implies
\[
\mu_p\left(\frac{\lambda}{1+\eps}\right)
\leq \rho(Fx,Fy)
\leq \mu_p((1+\eps)\lambda).
\]
Since \(0<\gamma\eps\leq1\), after reducing the concentration constant if
necessary, Lemma~\ref{lem:concentration} gives for this fixed pair
\[
\Prob\{\text{the desired inequality fails for }x,y\}
\leq 2\exp(-c_p\gamma^2 k\eps^2).
\]
There are fewer than \(N^2/2\) unordered pairs.  Choosing
\[
  k\geq C(p,\lambda_0,a)\eps^{-2}\log N
\]
with \(C\) large enough makes the union bound at most \(N^{-a}\).  Since
\((1-\eps)^{-2}\geq1\), the displayed dimension bound in the theorem is a
slightly weaker sufficient condition with the same dependence as Casey's
Theorem 1.0.1.  This proves the simultaneous assertion.
\end{proof}

\section{Verification and novelty notes}

The proof uses only standard p-stable facts: the scaling identity
\(\sum_j v_j W_j\stackrel{d}{=}\norm v_p W\), finite moments
\(\E\abs W^q<\infty\) for \(q<p\), and boundedness of the p-stable density near
zero.  These are the same stable-distribution background facts invoked by the
source paper and by standard references on p-stable projections.

Bounded novelty search performed before promotion:
\begin{itemize}
\item local run indexes were searched for arXiv:1906.03536, the exact title,
  ``p-stable'', ``metric preserving'', and ``Conjecture 1.0.2'';
\item the parsed source TeX and bibliography for arXiv:1906.03536 were
  inspected;
\item web/arXiv searches were run for exact-title and close-variant phrases,
  including ``p-stable metric preserving dimension reduction'', ``Casey
  p-stable conjecture'', and ``Cauchy projections metric preserving
  function''.
\end{itemize}
No exact later answer or duplicate packet was found in these bounded searches.

Human-review recommendation: review as a likely valid full solution to the
fixed-threshold reading of Conjecture 1.0.2.  The main points to check are the
uniform sub-exponential domination by \(\abs{\log\abs W}\), the gap estimate in
Lemma~\ref{lem:gap}, and the interpretation of \(\Omega(1)\).

\begin{thebibliography}{9}
\bibitem{Casey}
Michael P. Casey.
\newblock Linear dimension reduction approximately preserving a function of the
1-norm.
\newblock \emph{Electronic Journal of Statistics}; arXiv:1906.03536.

\bibitem{Indyk}
Piotr Indyk.
\newblock Stable distributions, pseudorandom generators, embeddings, and data
stream computation.
\newblock \emph{Journal of the ACM} 53 (2006), 307--323.

\bibitem{LedouxTalagrand}
Michel Ledoux and Michel Talagrand.
\newblock \emph{Probability in Banach Spaces}.
\newblock Springer, 1991.

\bibitem{Pisier}
Gilles Pisier.
\newblock On the dimension of the \(\ell_p^n\)-subspaces of Banach spaces, for
\(1\leq p<2\).
\newblock \emph{Transactions of the American Mathematical Society} 276 (1983),
201--211.
\end{thebibliography}

\end{document}
